\documentclass[11pt,a4paper]{article}
\usepackage[margin=2.4cm]{geometry}
\usepackage{amsmath,amssymb,booktabs,array,graphicx,xcolor}
\usepackage[colorlinks=true,linkcolor=blue!60!black,urlcolor=blue!60!black]{hyperref}
\hypersetup{pdftitle={Matching Toda edge scattering to good-Boussinesq semiclassics},
pdfauthor={Quantum Boussinesq solitons project}}
\newcommand{\Lc}{\textsf{\textbf{(c)}}}
\newcommand{\Ld}{\textsf{\textbf{(d)}}}
\newcommand{\scr}[1]{\texttt{matching\_#1.py}}
\newcommand{\Li}{\operatorname{Li}}
\newcolumntype{L}[1]{>{\raggedright\arraybackslash}p{#1}}
\emergencystretch=2em
\title{Matching the Toda edge phase to good-Boussinesq semiclassics:\\a nonuniform limit and a sign correction}
\author{Quantum Boussinesq solitons project: Phase 2f}
\date{3 November 2026}
\begin{document}
\maketitle
\begin{abstract}
The proposed quantitative matching does not close in the conventions currently
used by the project. The good-equation phonon phase has a finite small-amplitude
limit, but its complete one-loop coefficient at fixed physical momentum diverges
logarithmically in both collision channels. This follows from the existing
one-soliton momentum shift and the full Darboux shifts, checked independently by
differentiating the classical phase. At fixed edge variable $y=\alpha^2/\hbar$,
the relative momentum correction is itself unbounded, so taking the edge limit
before the semiclassical limit is not justified by this expansion. \Lc

We also correct a sign in the previously claimed cancellation to a pure
Legendre chi function. \Lc\ Independently, we scan the Toda Bethe phase at fixed
physical dressed momenta, retaining its existing classical frame conversion,
and separate mesh, edge-extrapolation and large-$y$ uncertainties. \Lc\ (numerical)
These tests
do not prove inequivalence of exact quantum theories, nor identify an exact
dressing factor. They locate the missing step: a controlled quantum momentum
and frame dictionary.~\Ld
\end{abstract}

\section{What is being compared}
All calculations here use the existing equations and normalisations; no new
regularisation prescription is chosen. The input is the good-equation
one-loop result in \texttt{oneloop\_phase.tex}, and the Toda computation in
\texttt{toda\_edge.tex}. The latter approaches the \emph{bad} equation.
Their identification at the quantum level is a hypothesis \Ld, not an input.

\paragraph{Classical dictionary (\Lc, \scr{conventions}).}
For the right-moving long-wave scaling
\[
 u=\lambda^2 U(X,T),\qquad X=\lambda(x-t),\qquad T=\lambda^3t,
\]
the good and bad equations, respectively, give
\[
 U_T=+\tfrac12 U_{XXX}+3UU_X,
 \qquad U_T=-\tfrac12 U_{XXX}-3UU_X.
\]
Here the integration constant is zero for decaying pulses. The leading
equations are related by reversing slow time, not by identifying colliding
particle labels. In the good equation the smaller pulse is faster; in the bad
equation the larger pulse is faster.

Write $0<r<1$ and take $(\alpha_1,\alpha_2)=(r\lambda,\lambda)$ for a
colliding good-equation pair, and $(\beta_1,\beta_2)=(\lambda,r\lambda)$ for
a colliding bad-equation pair. Define
\[
 A_0(r)=\left(\frac{1-r}{1+r}\right)^2.
\]
The full classical phases in the project then satisfy
\begin{align}
 \delta_{0,g}^{\rm same}&=\lambda^2(r^2-1)\log A_0+O(\lambda^4),\label{eq:classsame}\\
 \delta_{0,b}^{\rm same}&=\lambda^2(1-r^2)\log A_0+O(\lambda^4),\\
 \delta_{0,g}^{\rm head}&=4r\lambda^2+O(\lambda^4),\qquad
 \delta_{0,b}^{\rm head}=4r\lambda^2+O(\lambda^4).
\end{align}
The head-on ordering is $(+, -)$ in both equations. Thus even the classical
comparison needs a channel-dependent ordering dictionary. These equations
alone do not establish a quantum conjugation or continuation rule.

For Toda, retain the established dictionary
$\kappa=\sqrt3\,\beta$, $b=e^{-a/2}$, $\hbar_T=3b\hbar_B$,
and $y_j=\beta_j^2/\hbar_B$. All new runs use $a=1$ and determine excitations
from $P_{\rm dr}(p)-P_{\rm dr}(e)=P_n(\kappa)/a$.
Here $P_n=4s b(\kappa\cosh\kappa-\sinh\kappa)$.
The leading coordinate dictionary is $x_B=2\sqrt3\,n$,
$t_B=2\sqrt3\,b\,t_T$; correspondingly
$P_n=6\sqrt3\,b\,P_B[1+O(\beta^2)]$.
These charge and velocity limits are checked in \scr{conventions}; they
do not supply their quantum corrections.
The overtaking phase is the same two-term subtraction used in Phase 2d:
\begin{equation}
 \Psi_T=\Phi(p_1|p_2)-\Phi(p_1|e_2)
 -M_1\big(P_{n2}/a-\Pi_2\big),\qquad
 M_j=-2\kappa_j,\quad \Pi_j=2s_jb\sinh\kappa_j.
\end{equation}
This removes the first particle's edge-reference phase. The final term is
\emph{still the classical frame conversion}; its quantum completeness is not
established. Head-on scattering uses the existing four-term edge subtraction.

\section{The complete good-equation coefficient}
In this section \Lc\ means an algebraic consequence of the one-loop formulas
already adopted in Phase 2e, rather than a new independent derivation of the
fluctuation determinant. All new limits and signs are checked in
\scr{conventions} with 60-digit arithmetic.

The input physical-momentum convention is
\[
 P^{\rm phys}=P(\alpha)+\hbar\Delta P(\alpha),\qquad
 P(\alpha)=\tfrac23\epsilon\sin^3\alpha\cos\alpha.
\]
Put $c=-\log2-1+\tfrac14\log3$. The existing closed-form momentum shift gives
\begin{equation}
 \Delta P(\alpha,\epsilon)
 =\frac{\epsilon\alpha}{\pi}(\log\alpha+c)
 +O(\alpha^3\log\alpha).\label{eq:dp}
\end{equation}
The coefficient in $\arg S=\delta_0/\hbar+\delta_1+O(\hbar)$, when re-expressed
at fixed $P^{\rm phys}$, is
\begin{equation}
 \delta_1=I_{12}+\delta_{\rm ren},\qquad
 \delta_{\rm ren}=-\sum_j\frac{\partial\delta_0}{\partial P_j}\Delta P_j
 =\sum_j\Delta X_j'\Delta P_j.\label{eq:ren}
\end{equation}
The plus sign in the final expression is essential: the definition is
$\Delta X_j'=-\partial\delta_0/\partial P_j$.
Direct differentiation agrees with the implemented Darboux formula to
better than $10^{-50}$ in the new test.

\subsection{Overtaking: a term lost by keeping only the Hirota shift}
For $(\alpha_1,\alpha_2)=(r\lambda,\lambda)$, the Darboux coefficient
$\gamma_j=-3\cos2\alpha_j/(\sin\alpha_j\sin3\alpha_j)$ obeys
$\gamma_j=-\alpha_j^{-2}+O(1)$. Separate \eqref{eq:ren} into the Hirota
shift and the remaining Darboux contribution:
\begin{align}
 H(r)&=-\frac{\log A_0(r)\log r}{\pi},\label{eq:hirota}\\
 D(r,\lambda)&=\frac2\pi\left[
 (r-r^{-1})(\log\lambda+c)-\frac{\log r}{r}\right].\label{eq:darboux}
\end{align}
The raw phonon term tends to $I_0(r)=2L(r)/\pi$, where
\[
 L(r)=-2\chi_2(r)-\log r\log\frac{1-r}{1+r},\qquad
 \chi_2(r)=\tfrac12\big[\Li_2(r)-\Li_2(-r)\big].
\]
Consequently the \emph{complete} coefficient is
\begin{equation}
 \boxed{\delta_1^{\rm same}(r\lambda,\lambda)
 =\frac2\pi L(r)+H(r)+D(r,\lambda)
 +O(\lambda^2|\log\lambda|).}\label{eq:samefull}
\end{equation}
In particular, its logarithmic slope is
$\partial_{\log\lambda}\delta_1\to2(r-r^{-1})/\pi$.
It depends on the amplitude ratio and cannot be removed by a constant phase.
Replacing $I$ by the antisymmetrised $J=I+\pi/2$ does not change this conclusion.

\paragraph{Correction to Phase 2e.}
The paragraph following its small-amplitude $\Delta P$ formula claims that the
Hirota correction cancels the logarithmic part of $I_0$, leaving
$-4\chi_2(r)/\pi$. With its stated signs, the actual sum is
\begin{equation}
 \boxed{I_0(r)+H(r)=-\frac4\pi\chi_2(r)
 -\frac4\pi\log r\log\frac{1-r}{1+r}.}\label{eq:signcorrection}
\end{equation}
At $r=1/2$ the actual value is $-1.62570710834$, whereas the claimed pure-chi
value is $-0.65613518176$. The old limits script tests $H$ and $I$ separately
but never asserts their claimed cancellation. Its passing status therefore
does not validate that sentence. The original report is left intact here;
this note and the README explicitly record the correction.

\subsection{Head-on: the phonon limit is not the physical-momentum limit}
For $(\alpha_1,\epsilon_1)=(r\lambda,+1)$ and
$(\alpha_2,\epsilon_2)=(\lambda,-1)$, the phonon term tends to zero, but
\begin{equation}
 \boxed{\delta_1^{\rm head}
 =-\frac2\pi\left[(r+r^{-1})(\log\lambda+c)+\frac{\log r}{r}\right]
 +O(\lambda^2|\log\lambda|).}\label{eq:headfull}
\end{equation}
Thus the vanishing fixed-$\alpha$ phonon correction cited in the README is not
by itself a match to a fixed-physical-momentum Toda phase.
Both \eqref{eq:samefull} and \eqref{eq:headfull}, including their ratio
derivatives, were checked at $r=1/4,1/2,3/4$ and
$\lambda=10^{-3},10^{-6},10^{-9}$.
For example, at $r=1/2$ the full expressions give:
\begin{center}\small
\begin{tabular}{@{}crrr@{}}\toprule
$\lambda$ & $\delta_1^{\rm same}$ & $\delta_1^{\rm head}$ &
Largest asymptotic error\\\midrule
$10^{-3}$ & $7.20781548$ & $14.13419003$ & $1.02\times10^{-5}$\\
$10^{-6}$ & $13.80423817$ & $25.12821386$ & $1.96\times10^{-11}$\\
$10^{-9}$ & $20.40065856$ & $36.12224784$ & $2.90\times10^{-17}$\\\bottomrule
\end{tabular}
\end{center}

\subsection{Why the proposed order of limits fails}
From \eqref{eq:dp}, with $y=\alpha^2/\hbar$,
\begin{equation}
 \frac{\hbar\Delta P}{P}
 =\frac{3}{2\pi y}(\log\alpha+c)
 +O\!\left(\frac{\alpha^2|\log\alpha|}{y}\right).\label{eq:nonuniform}
\end{equation}
At fixed finite $y$ this ratio is unbounded as $\alpha\to0$. A necessary
condition for this perturbative momentum conversion is instead
$y\gg|\log\alpha|$. The nested limit ``edge first at fixed $y$, then
$y\to\infty$'' leaves its domain of validity before the second step.
Even the zero-amplitude edge subtraction of the full good-equation coefficient
requires care, since \eqref{eq:darboux} is singular at that endpoint.

This is a failure of a \emph{termwise comparison}, not a prediction that an
exact physical scattering amplitude diverges. A resummation or a quantum
dictionary could change the conclusion; none is selected or assumed here. \Ld

\section{Independent Toda scan}
\scr{toda} retains the existing digamma-kernel solver and computes
\begin{equation}
 R_h(Y,r)=\frac{\Psi_T}{\hbar_T}
 -Y(1-r^2)\log A_0(r),\qquad (y_1,y_2)=(Y,r^2Y).
\end{equation}
This definition subtracts the limiting classical phase, not a fitted
polynomial or a candidate quantum amplitude. A fit in $\hbar_T$ at fixed
$(Y,r)$ removes finite-amplitude corrections before examining large $Y$.
The values $Y=4,8,16,32$ and $r=0.4,0.5,0.6$ keep the first comparison
away from both the equal-velocity and zero-amplitude endpoints.

\paragraph{Resolution and order of limits.}
The edge sequence is $\hbar_T=0.04,0.02,0.01,0.005,0.0025$,
with 681, 1062, 1824, 3348, 6397 quadrature nodes, respectively.
At $\hbar_T=0.01$, increasing the node count from 1824 to 2554 changes all sampled phases by at most $2.3\times10^{-12}$ radians. This mesh test is at one fixed $\hbar_T$;
it is not a convergence guarantee at every smaller $\hbar_T$.
The edge sequence uses the original NumPy Gauss--Legendre nodes.
At $\hbar_T=0.01$ and 1824 nodes, the NumPy/SciPy phase difference is $1\times10^{-12}$ radians.
The central edge estimate is the quadratic intercept from the three
smallest $\hbar_T$. The stated fit spread is the largest change on using
the preceding quadratic window, the four-smallest-point cubic fit, or
a fit to $1,\hbar_T,\hbar_T\log\hbar_T$ at the three smallest points.
The last-two-point linear fit is stored separately in the JSON diagnostics.
\begin{center}\small
\begin{tabular}{@{}ccrr@{}}\toprule
$Y$ & $r$ & Edge residual $R_0$ & Edge fit spread\\\midrule
4 & 0.4 & -0.011162 & $3.0\times10^{-6}$\\
4 & 0.5 & -0.024831 & $5.0\times10^{-6}$\\
4 & 0.6 & -0.048097 & $6.8\times10^{-6}$\\
8 & 0.4 & -0.010497 & $5.1\times10^{-6}$\\
8 & 0.5 & -0.022191 & $4.6\times10^{-6}$\\
8 & 0.6 & -0.041615 & $3.4\times10^{-6}$\\
16 & 0.4 & -0.010090 & $3.2\times10^{-5}$\\
16 & 0.5 & -0.020857 & $9.8\times10^{-5}$\\
16 & 0.6 & -0.038509 & $2.3\times10^{-4}$\\
32 & 0.4 & -0.009868 & $5.4\times10^{-4}$\\
32 & 0.5 & -0.020201 & $1.3\times10^{-3}$\\
32 & 0.6 & -0.037002 & $2.7\times10^{-3}$\\
\bottomrule\end{tabular}\end{center}
\paragraph{Large-$Y$ diagnostic, not an exact limit.}
A quadratic fit in $1/Y$ to $Y=8,16,32$ gives the following exploratory
intercepts. The model spread compares a linear fit to the largest two
$Y$ values and a cubic fit to all four; edge-extrapolation uncertainty
is propagated with the absolute interpolation weights and shown separately.
\begin{center}\small\begin{tabular}{@{}crrr@{}}\toprule
$r$ & Exploratory $Y\to\infty$ intercept & Model spread & Edge spread\\\midrule
0.4 & -0.00964 & $1.2\times10^{-5}$ & $1.5\times10^{-3}$\\
0.5 & -0.01955 & $7.5\times10^{-6}$ & $3.7\times10^{-3}$\\
0.6 & -0.03553 & $3.1\times10^{-5}$ & $7.6\times10^{-3}$\\
\bottomrule\end{tabular}\end{center}
\paragraph{Derivative and head-on controls.}
At $Y=16,r=0.5$, the central ratio derivative is $-0.13765$ with step $0.01$, versus $-0.14210$ with step $0.1$.
The propagated edge-fit spread for the finer derivative is $9.8\times10^{-3}$.
This supports nonconstant ratio dependence independently of an additive phase.
The head-on residuals after subtracting $4\sqrt{y_1y_2}$ are:
\begin{center}\small\begin{tabular}{@{}crr@{}}\toprule
$(y_1,y_2)$ & Edge residual & Edge fit spread\\\midrule
$(1,1)$ & $-1.54\times10^{-8}$ & $8.4\times10^{-6}$\\
$(4,1)$ & $-1.61\times10^{-7}$ & $1.5\times10^{-5}$\\
$(16,4)$ & $-3.23\times10^{-5}$ & $1.3\times10^{-3}$\\
\bottomrule\end{tabular}\end{center}
These controls support the previous numerical head-on result in the
Toda extraction convention. They do not establish the full good-equation match.


\paragraph{Interpretation (\Lc\ numerical / \Ld\ identification).}
The tabulated residual is a property of the existing Toda phase extraction,
including its classical frame conversion. A finite, ratio-dependent residual
is numerical evidence, not an analytical determination of its large-$Y$
limit. Neither adding a constant nor changing the overall overtaking sign
can repair the nonuniform physical-momentum expansion in
\eqref{eq:nonuniform}. We therefore do not fit the Toda residual to a new
quantum dilogarithm or claim that it confirms the good-equation phonon phase.

\section{Outcome, checks, and next dependency}
The proposed matching test has a reproducible negative outcome: the current
formulas do not supply a controlled common quantum limit. This conclusion is
more specific than a disagreement between two plotted curves. It identifies
the divergent momentum conversion, preserves the distinction between the two
equations, and corrects a definite algebraic sign error. \Lc

The next dependency is to derive the quantum label-frame and momentum
dictionary in a common scaling regime, checking whether the logarithmic terms
are absorbed, survive, or signal inequivalent quantum limits. That question
must precede using the Toda residual to constrain a finite-amplitude exact
$S$-matrix. A joint limit satisfying $y\gg|\log\alpha|$ is a possible test,
but is not the nested limit previously proposed. \Ld

\begin{center}\small
\begin{tabular}{@{}L{6.3cm}L{2cm}L{4.7cm}@{}}\toprule
Claim & Label & Check\\\midrule
Classical ordering and slow-time dictionary & (c) & \scr{conventions}\\
Full logarithmic coefficients and ratio derivatives & (c) & \scr{conventions}, 60 digits\\
Correction to pure-chi cancellation & (c) & \scr{conventions}\\
Nonuniform momentum expansion & (c) & \scr{conventions}\\
Toda edge residuals and error diagnostics & (c), numerical & \scr{toda}; \scr{analysis}\\
Existence of a common exact quantum limit & (d), unresolved & Not assumed\\\bottomrule
\end{tabular}
\end{center}

\paragraph{Failures and limitations recorded.}
\begin{itemize}
\item The expected direct matching fails its convention and uniformity check:
the full physical-momentum coefficient has no finite fixed-$y$ edge limit.
\item The pure-chi cancellation asserted in Phase 2e fails by about $0.970$
at $r=1/2$. The new regression check verifies the corrected expression;
it does not weaken the old checks or change their signs.
\item Mesh refinement, different edge fits and different large-$Y$ fits
measure different errors. Fit spreads below are diagnostics, not rigorous
confidence intervals, and do not resolve an unknown quantum frame term.
\item The original 6397-node run took about 27 minutes. A redundant fine-grid
run with alternative node construction was stopped once that result was
available; no unfinished fine-grid result is used. The alternative method
was separately completed at $\hbar_T=0.01$ for the node-method comparison.
\item The good-equation formulas are accepted as inputs. Their algebra and
limits are tested here, but the FK determinant derivation is not re-proved.
\end{itemize}

\clearpage
\section*{Reproduction and provenance}
The saved JSON files allow the analytic checks and all extrapolations to be
reproduced without repeating the expensive dense solves. From the repository
root, run:
\begin{verbatim}
cd phase2/scripts
python3 matching_conventions.py
python3 matching_analysis.py
\end{verbatim}
To regenerate the raw Toda data as well, run the following from
\texttt{phase2/scripts/}. The node-method check and mesh refinement are
independent computations; the fine-grid run can take tens of minutes.
\begin{verbatim}
export OPENBLAS_NUM_THREADS=1 OMP_NUM_THREADS=1
python3 matching_toda.py
python3 matching_toda.py --hs .0025 \
  --output matching_toda_fine.json
python3 matching_toda.py --hs .01 --factor 1.4 \
  --output matching_toda_nodes.json
python3 matching_toda.py --hs .01 --fast-nodes \
  --output matching_toda_fast.json
python3 matching_analysis.py
\end{verbatim}
Each completed checking script ends in \texttt{ALL CHECKS PASSED}.
The existing \path{oneloop_good_twosoliton_phase.py} and
\path{oneloop_good_limits_cusp_resonance.py} were also rerun successfully.
The analytic checks use 60-digit \texttt{mpmath}; the dense Bethe solves use
double-precision NumPy/SciPy. Raw data retain node counts, ground-state
normalisation residuals, phases, and elapsed times. Derived JSON records the
alternative fits as well as the central estimates.

To build the report after regenerating the tables:
\begin{verbatim}
cd ..
pdflatex matching.tex
pdflatex matching.tex
\end{verbatim}

The sources inspected for this work are the project reports
\texttt{report.tex}, \texttt{toda\_route.tex}, \texttt{toda\_edge.tex},
\texttt{oneloop\_phase.tex}, and their relevant scripts. No new external
literature result or bibliographic reference is introduced.
\end{document}
